% Part: history
% Chapter: biographies 
% Section: georg-cantor

\documentclass[../../../include/open-logic-section]{subfiles}

\begin{document}

\olfileid{his}{bio}{can}

\olsection{గెయోర్గ్ కాంటర్}

\olphoto{cantor-georg}{గెయోర్గ్ కాంటర్}

గెయోర్గ్ కాంటర్ (ఉచ్చారణ: గే-ఒర్గ్ కాన్-టోర్) గురించి వచ్చిన తొలి జీవిత చరిత్ర ఒకటి, ఆయనను రష్యాలోని సెయింట్ పీటర్స్‌బర్గ్‌కు వెళ్తున్న ఓడలో పుట్టిన శిశువుగా కనుగొన్నారని, ఆయన తల్లిదండ్రులెవరో తెలియరని చెప్పింది. కానీ అది నిజం కాదు; అయితే ఆయన 1845లో సెయింట్ పీటర్స్‌బర్గ్‌లో పుట్టారు.

కాంటర్ 1867లో బెర్లిన్ విశ్వవిద్యాలయంలో గణితంలో డాక్టరేట్ పొందారు. సమితి సిద్ధాంతంలో ఆయన చేసిన కృషికి ప్రసిద్ధి; సమితి సిద్ధాంతాన్ని ప్రత్యేక పరిశోధనా విభాగంగా స్థాపించిన ఘనత ఆయనదిగా భావిస్తారు. వేర్వేరు పరిమాణాల అనంత సమితులు ఉంటాయని తొలిసారిగా నిరూపించింది ఆయనే. ఆయన సిద్ధాంతాలు, ముఖ్యంగా అనంతాల గురించిన సిద్ధాంతం, అప్పటి గణితశాస్త్రజ్ఞులలో విస్తృత చర్చకు దారితీశాయి; ఆయన కృషి వివాదాస్పదమైంది.

కాంటర్ మత విశ్వాసాలు, గణిత కృషి విడదీయలేనంతగా ముడిపడ్డాయి; అతిపరిమిత సంఖ్యల సిద్ధాంతాన్ని దేవుడే తనకు నేరుగా తెలియజేశాడని కూడా ఆయన చెప్పారు. జీవితపు తరువాతి భాగంలో కాంటర్ మానసిక అనారోగ్యంతో బాధపడ్డారు. 1894 నుంచి, మరీ ముఖ్యంగా చివరి సంవత్సరాలలో, ఆయన పలుమార్లు ఆసుపత్రిలో చేరారు. గణితశాస్త్రజ్ఞుడు లియోపోల్డ్ క్రోనెకర్‌తో విభేదంతో సహా, ఆయన కృషిపై వచ్చిన తీవ్రమైన విమర్శలు నిస్పృహకు, గణితంపై ఆసక్తి తగ్గడానికి దారితీశాయి. నిస్పృహ కాలాల్లో కాంటర్ తత్వశాస్త్రం, సాహిత్యం వైపు మళ్లారు; షేక్స్‌పియర్ నాటకాలను నిజానికి ఫ్రాన్సిస్ బేకన్ రాశాడనే సిద్ధాంతాన్ని కూడా ప్రచురించారు.

కాంటర్ 1918 జనవరి 6న హాలేలోని ఒక చికిత్సా కేంద్రంలో మరణించారు.

\begin{reading}
కాంటర్ సమగ్ర జీవిత చరిత్రలకు \citet{Dauben1990}, \citet{Grattan-Guinness1971} చూడండి. ఆయన విప్లవాత్మక అభిప్రాయాలను బీబీసీ రేడియో 4 కార్యక్రమం \emph{A Brief History of
  Mathematics}లో కూడా వివరించారు \citep{Sautoy2014}. కాంటర్ సిద్ధాంతాలను ర్యాప్ రూపంలో వినాలనుకుంటే \citet{Rose2012} చూడండి.
\end{reading}

\end{document}
